\chapter{Come usare questo libro}

Non si presuppone alcuna conoscenza. Chi non ha mai programmato, o non ha mai
guardato dentro un computer, cominci dal capitolo \emph{Partire da zero}, a
pagina~\pageref{ch:zero}: poche pagine che spiegano le parole e le idee usate
fin dalla prima lezione.

\section*{Un capitolo, una lezione}

Il libro ha venticinque capitoli, numerati da~0 a~24 come le lezioni del corso,
e come quasi tutto il resto in questo libro: bit, registri, righe e colonne
partono tutti da zero. Il Capitolo~$N$ segue la Lezione~$N$, che ha due
presentazioni. La \emph{presentazione della lezione}, di al più venti slide,
contiene i concetti; la \emph{presentazione del codice} contiene tutto ciò che
si scrive e si esegue: ogni listato con la sua spiegazione, il programma
completo della lezione, che cosa mostra, e il modo di provarlo su \prgqt, su
QEMU e sulla scheda. Il capitolo è la forma estesa di entrambe: l'ordine dei
paragrafi è l'ordine delle slide, i listati sono gli stessi file, le domande
alla fine sono quelle mostrate in aula, e ciò per cui in venti slide non c'è
spazio è qui. Chi ha seguito la lezione trova le frasi che una slide non può
contenere; chi l'ha persa trova la lezione.

La maggior parte dei capitoli termina con un paragrafo, \emph{Il programma di
questo capitolo}, che stampa per intero il programma da cui sono tratti i
listati del capitolo, mostra il frame che esso produce e dà i comandi per
eseguirlo. I programmi sono nella cartella \fname{code} del repository del
corso, uno per lezione, e ciascuno è compilato ed eseguito dai test automatici;
il Paragrafo~\ref{sec:prg:testing} è la guida passo per passo per i tre host.

\begin{center}\small
\begin{tabularx}{\textwidth}{@{}lP{0.36\textwidth}L@{}}
\toprule
\textbf{Parte} & \textbf{Capitoli} & \textbf{Poi, il progetto}\\
\midrule
Partenza & 0 Introduzione, 1 Installazione, 2 Primo contatto & \\
Informazione & 3 Storia, 4 Informazione digitale, 5 Interi, 6 Reali & 7 \vermone, parte 1\\
Logica & 8--9 Algebra di Boole, 10 Elettronica, 11 Macchine combinatorie & 12 \vermone, parte 2\\
Stato e memoria & 13 Macchine sequenziali, 14 Registri e bus, 15 Memoria, 16 Memoria permanente & 17 \vermone, parte 3\\
La macchina & 18 Processore, 19 Istruzioni e interruzioni, 20 Scheda video, 21 Ingresso e uscita & 22 \vermone, parte 4\\
Pubblicazione, prospettive & 24 Oltre un solo core & 23 \vermone sul Cartridge Store\\
\bottomrule
\end{tabularx}
\end{center}

\section*{La forma di un capitolo di teoria}

Ogni capitolo che presenta un argomento ha gli stessi cinque movimenti.

\begin{enumerate}
\item \textbf{Riscaldamento sulla macchina.} Poche righe di \riscv su \prg che
  producono un fatto che chiede una spiegazione. Eseguitele prima di
  proseguire.
\item \textbf{La trattazione formale.} Definizioni, teoremi, metodi ed esempi
  svolti.
\item \textbf{Ritorno alla macchina.} Le stesse regole, come le realizza
  \riscv.
\item \textbf{Verificate di aver capito.} Domande a risposta multipla nello
  stile della prova intercorso. Le risposte, ciascuna con una riga di
  spiegazione, sono nell'Appendice~\ref{app:answers}: leggetele dopo aver
  risposto, non al posto di rispondere.
\item \textbf{Esercizi, glossario, letture.} Gli esercizi si svolgono prima su
  carta e poi su \prg; il glossario affianca a ogni termine inglese quello
  italiano, perché la documentazione tecnica, i messaggi degli strumenti e i
  commenti dei listati sono in inglese.
\end{enumerate}

\section*{La forma di un capitolo tutorial}

I Capitoli 7, 12, 17, 22 e 23 costruiscono \vermone. Sono autonomi e pensati
per essere seguiti alla tastiera. Ogni passo mostra poche righe, dice se sono
\textbf{nuove}, \textbf{modificate} o \textbf{invariate}, e finisce con un
punto di controllo \emph{Eseguite} che dichiara che cosa deve mostrare lo
schermo. Se all'inizio di un tutorial un programma si rifiuta di funzionare,
ripartite dal file di riferimento della parte precedente
(Appendice~\ref{app:verm1}): nessuno deve restare indietro per un errore.

\section*{I riquadri}

\begin{prgconcept}
Un riquadro \textbf{Concetto} enuncia un'idea da cui dipende il resto del
capitolo.
\end{prgconcept}
\begin{definition}{come appaiono le definizioni}
Definizioni formali e teoremi sono numerati all'interno del capitolo.
\end{definition}
\begin{prgtry}
Un riquadro \textbf{Provate voi} chiede di fare qualcosa alla tastiera, e dice
che cosa si deve osservare.
\end{prgtry}
\begin{prgwarn}
Un riquadro \textbf{Attenzione} segnala un errore che si ripete ogni anno.
\end{prgwarn}
\begin{prgmachine}
Un riquadro \textbf{Ritorno alla macchina} lega un risultato formale a
un'istruzione, a un registro o a una schermata del debugger.
\end{prgmachine}
\begin{prgaside}
Un riquadro \textbf{Nota storica} racconta da dove viene un'idea.
\end{prgaside}

\section*{Stato di apprendimento: comprendere oggi, utilizzare oggi, comprendere più avanti}

Il corso prima programma e poi formalizza, quindi una lezione mostra sempre più di quanto
spieghi. Per questo ogni capitolo si apre, dopo gli obiettivi, con un riquadro che divide ciò
che contiene in tre stati, e le slide portano le stesse etichette sopra il titolo.

\begin{description}[style=nextline,leftmargin=1.5em]
\item[\masterychip{now}{DA COMPRENDERE OGGI}] Alla fine della lezione sapete spiegarlo,
  seguirlo a mano passo passo e applicarlo da soli. Le domande e gli esercizi verificano questi.
\item[\masterychip{use}{DA UTILIZZARE OGGI}] Oggi seguite e usate la ricetta: una direttiva, un
  comando, una chiamata del runtime, una convenzione sui registri. Come funziona dentro lo
  spiega il capitolo indicato tra parentesi, oppure non è oggetto del corso (gli strumenti di
  \prg si usano, non si studiano).
\item[\masterychip{later}{DA COMPRENDERE PIÙ AVANTI}] Un'anticipazione: riconoscete il nome e lo
  scopo. Lo spiega il capitolo indicato tra parentesi; nulla è rinviato senza una destinazione.
\end{description}

Gli stati sono attese nel tempo, non livelli di difficoltà: la stessa idea passa spesso da
\emph{più avanti} a \emph{utilizzare} a \emph{comprendere} nel giro di qualche capitolo (la
convenzione di chiamata si usa dal Capitolo~\ref{ch:firstcontact} e si studia nel
Capitolo~\ref{ch:processor}). Una slide o un paragrafo senza etichetta è contesto: una storia,
un esempio, una domanda, un esercizio. Il colore non è mai l'unico indizio: il libro, la chiave
inglese e l'orologio, e le cornici continua, doppia e tratteggiata, distinguono i tre stati
anche in una stampa in bianco e nero.

\section*{Notazione}

I numeri senza indicazioni sono decimali, con la virgola decimale e, dove
serve, il punto per le migliaia. I numeri esadecimali sono scritti come li
vuole l'assemblatore, \hex{07E0}; i numeri binari come \code{0b1010} nel codice
e con un pedice, $1010_2$, nelle formule. I bit si numerano da destra partendo
da zero, così il bit~$i$ ha peso $2^{i}$. Registri (\reg{a0}, \reg{sp}),
istruzioni (\ins{addi}) e nomi di file sono composti in carattere a spaziatura
fissa. Nelle espressioni booleane il complemento di $x$ è $\nott{x}$, la somma
è $+$ e il prodotto si scrive per giustapposizione. I multipli del byte seguono
i prefissi binari dove la precisione conta:
$1~\mathrm{KiB}=2^{10}$ byte, $1~\mathrm{MiB}=2^{20}$ byte.

I listati e i loro commenti sono in inglese, come gli identificatori di \prg,
i comandi e le etichette dell'interfaccia di \prgqt (\key{Step}, \key{Resume},
\key{Memory}): sono ciò che compare davvero sullo schermo.

\section*{Le appendici}

L'Appendice~\ref{app:riscv} è un riferimento completo dell'insieme di
istruzioni \riscv usato nel libro: registri, formati, ogni istruzione di RV32I
e delle estensioni M, A, C e Zicsr, pseudo-istruzioni, direttive
dell'assemblatore, la convenzione di chiamata e i registri di controllo del
modo macchina. L'Appendice~\ref{app:prg32} è il riferimento di \prg: il
contratto, l'interfaccia di programmazione, il formato delle cartucce, gli
strumenti a riga di comando, gli host e il Cartridge Store.
L'Appendice~\ref{app:github} è un riferimento di Git e GitHub, gli strumenti
con cui la piattaforma, questo libro e il progetto del lettore sono conservati
e condivisi. L'Appendice~\ref{app:verm1} stampa per intero i quattro programmi
di riferimento di \vermone; l'Appendice~\ref{app:answers} le risposte a tutte
le domande; l'Appendice~\ref{app:glossary} il glossario completo
inglese--italiano.
